\documentclass[coursework]{SCWorks}
\usepackage{preamble}

\begin{document}

% ---------- Титульный лист ----------
\chair{математической кибернетики и компьютерных наук}
\title{Анализ аварийности на дорогах Саратова по структурным признакам дорожной сети}
\course{3}
\group{341}
\napravlenie{02.03.03 "--- Математическое обеспечение и администрирование информационных систем}
\author{Колугова Егора Алексеевича}

% Заведующий кафедрой
\chtitle{доцент, к.\,ф.-м.\,н.}
\chname{М.\,В.\,Огнева}

% Научный руководитель — должность/степень и инициалы с фамилией
\satitle{ст. преподаватель}
\saname{А.\,А.\,Казачкова}

\date{2027}

\maketitle
\secNumbering
\tableofcontents

% ---------- Введение ----------
\intro

Ежегодно в России в дорожно-транспортных происшествиях погибает более
десяти тысяч человек~\cite{gibddstat}. Значительная часть ДТП
концентрируется на небольшом числе участков улично-дорожной сети.
Для выявления таких участков и выбора мер по их переустройству
используются отечественные нормативные методики и зарубежные
руководства, в том числе Highway Safety Manual.

Классические методики предусматривают как обследование отдельных
объектов, так и автоматизированный скрининг дорожной сети. Однако для их
применения часто нужны подробные характеристики дорог и интенсивности
движения, подготовка которых трудоёмка. Открытые источники позволяют
автоматизировать часть этой работы: OpenStreetMap описывает геометрию и
атрибуты сети, федеральная информационная система СКДФ содержит
паспортные сведения, цифровые модели поверхности позволяют оценивать
рельеф, а карточки ДТП содержат координаты происшествий. Полнота этих
данных неодинакова между территориями. Это ставит вопрос о качестве
модели при обучении и применении в разных городах.

\textbf{Объект исследования} "--- аварийность на улично-дорожной сети
города Саратова. \textbf{Предмет исследования} "--- зависимость
пространственного распределения ДТП от структурных характеристик
дорожной сети, извлекаемых из открытых геоданных.

\textbf{Цель работы} "--- построить и обосновать модель, которая по
структурным признакам дорожной сети ранжирует ячейки территории города
по ожидаемому числу ДТП, оценить вклад групп признаков в качество
ранжирования и проверить переносимость подхода на другие города.
Межгородское сравнение служит внешней проверкой саратовской модели.
Термин «риск» далее обозначает ожидаемую частоту ДТП с пострадавшими
за заданный период.

Для достижения цели в теоретической части решаются следующие
\textbf{задачи}:
\begin{enumerate}
  \item рассмотреть аварийность как объект анализа данных и выделить
        её статистические свойства, определяющие выбор методов;
  \item изучить отечественные и зарубежные методы оценки опасности
        участков дорог и сопоставить их предпосылки;
  \item обосновать выбор пространственной единицы анализа и
        рассмотреть открытые источники данных о дорожной сети;
  \item рассмотреть принципы отбора структурных признаков и проблему
        неизвестной экспозиции;
  \item рассмотреть методы машинного обучения, применимые к редким
        пространственным событиям, и способы оценки качества
        ранжирования;
  \item разграничить проверку будущих периодов на известной территории
        и проверку переноса на новый город;
  \item сформулировать постановку практической задачи и ограничения
        интерпретации её результатов.
\end{enumerate}

Работа состоит из введения, двух разделов теоретической части,
разделов практической части, заключения и списка использованных
источников. В первом разделе даётся обзор предметной области и
источников; во втором формулируется практическая задача, решение которой
описывается в последующих разделах.

% ---------- 1 Обзор предметной области ----------
\section{Обзор предметной области}

\subsection{Аварийность как объект анализа данных}

Дорожно-транспортное происшествие (ДТП) определяется законодательно как
событие, возникшее в процессе движения по дороге транспортного средства и
с его участием, при котором погибли или ранены люди, повреждены
транспортные средства, сооружения, грузы либо причинён иной материальный
ущерб~\cite{fz196}. Карточки учёта ДТП с пострадавшими содержат координаты,
дату и время, вид происшествия, сведения о дорожных условиях,
участниках и транспортных средствах~\cite{odm2015}. По данным
Госавтоинспекции, в 2023~году в
стране произошло около 132~тысяч учётных ДТП, в которых погибли более
14~тысяч и были ранены более 166~тысяч человек~\cite{gibddstat}.

В дорожной науке ДТП рассматривают как результат взаимодействия
человека, транспортного средства и дорожной среды~\cite{babkov1993,hsm2010}.
К дорожной среде относятся геометрия, покрытие и устройство пересечений.
Эти характеристики изменяются при строительстве, содержании
и переустройстве дорог.

ДТП "--- \emph{редкие события}. Если рассматривать
территорию, разбитую на небольшие участки или ячейки, и достаточно
короткие интервалы времени (например, месяц), то подавляющее большинство
пар «место "--- период» не содержит ни одного происшествия. Доля
положительных наблюдений в такой панели составляет единицы процентов или
доли процента. Доля правильных ответов при фиксированном пороге
малоинформативна: при доле положительных наблюдений 2"--~3\,\%
модель, всегда предсказывающая отсутствие ДТП, даст 97"--~98\,\%
правильных ответов, не обнаружив ни одного опасного места.

Число ДТП "--- \emph{счётная величина} с характерной
статистикой. Простейшей моделью числа событий на участке является
распределение Пуассона. При объединении неоднородных участков часто
наблюдается \emph{сверхдисперсия}: дисперсия числа ДТП превышает
среднее. Смесь пуассоновских распределений с разными интенсивностями
может объяснять такое свойство даже при пуассоновском процессе внутри
каждого участка. Поэтому следует различать разброс между местами и
остаточный разброс при условии известного уровня места. Для описания
неучтённой неоднородности часто используют отрицательное биномиальное
распределение, в частности в функциях безопасности (SPF). Выбор модели
должен учитывать остаточную дисперсию, а не только распределение всех
наблюдений вместе; обзор альтернатив дан в работе~\cite{lord2010}.

Аварийность \emph{зависит от экспозиции} "--- от объёма
движения. Чем больше автомобилей проходит через участок,
тем больше на нём происшествий при прочих равных условиях, причём
зависимость, как правило, нелинейна и растёт медленнее, чем
интенсивность движения, хотя форма зависимости различается между
типами объектов. Абсолютное число ДТП отражает одновременно условия
движения и загруженность. Относительные показатели нормируют его на
экспозицию, тогда как показатели частоты и критерии концентрации могут
использовать абсолютные числа.

ДТП также \emph{пространственно сконцентрированы}. Небольшая
доля участков сети собирает непропорционально большую долю
происшествий; именно это свойство лежит в основе понятия «место
концентрации ДТП» (очаг аварийности) и всех процедур выявления таких
мест~\cite{odm2015}. При этом концентрация в значительной степени
устойчива во времени: участок, опасный в прошлом году, с высокой
вероятностью останется опасным и в следующем, если на нём ничего не
изменилось.

Наблюдаемые числа ДТП подвержены \emph{случайным колебаниям}
и, как следствие, \emph{регрессии к среднему}. Участок, на котором в
отчётном году случайно произошло много ДТП, в среднем покажет
снижение в последующие годы даже без мероприятий; если такой участок был
отобран для реконструкции именно по высокому числу ДТП, то последующее
снижение будет ошибочно приписано реконструкции. Для оценки таких
изменений с поправкой на регрессию к среднему применяется
эмпирический байесовский метод~\cite{hauer1997}.

При работе с карточками ДТП существенны \emph{неполнота и неточность} данных.
Используемая выборка не охватывает происшествия без пострадавших;
координаты в карточках учёта определяются с ограниченной
точностью и иногда привязываются к ближайшему адресу или километру
дороги, а не к фактическому месту; часть полей заполняется
формально. Любая модель аварийности обучается на этой искажённой
выборке и наследует её искажения.

\subsection{Отечественные методы оценки опасности участков дорог}

Отечественная школа оценки безопасности дорог сложилась в 1960"--~1980-х
годах в работах В.\,Ф.~Бабкова и его последователей в МАДИ и была
закреплена в ведомственных строительных нормах ВСН~25-86 «Указания по
обеспечению безопасности движения на автомобильных
дорогах»~\cite{babkov1993,vsn2586}. В 2002~году на смену ВСН~25-86
пришли «Рекомендации по обеспечению безопасности движения на
автомобильных дорогах», утверждённые распоряжением Минтранса России
№~ОС-557-р~\cite{mintrans2002}; процедуры учёта и анализа самих ДТП
регламентирует более поздний отраслевой дорожный методический документ
ОДМ~218.6.015-2015~\cite{odm2015}.

Рекомендации Минтранса (п.~1.1.4) выделяют четыре метода выявления
опасных участков: метод на основе анализа данных о ДТП, метод
коэффициентов аварийности, метод коэффициентов безопасности и метод
конфликтных ситуаций. Первый применим к существующим дорогам при наличии
полной статистики происшествий за период не менее 3"--~5~лет; остальные
три предназначены для случаев, когда таких данных нет, а также для
оценки проектных решений.

\textbf{Метод на основе данных о ДТП.} Основной показатель "---
коэффициент относительной аварийности~$\mathit{И}$, число ДТП на один
миллион автомобиле"=километров пробега. Для протяжённых и однородных по
геометрическим элементам участков он вычисляется по формуле
(п.~1.1.5, формула~(1.1) Рекомендаций):
\begin{equation}
  \mathit{И} = \frac{10^{6}\, z}{365 \cdot L \cdot N},
  \label{eq:rel_accident_rate}
\end{equation}
где $z$ "--- число происшествий в год, $L$ "--- длина участка в
километрах, $N$ "--- среднегодовая суточная интенсивность движения в
обоих направлениях, авт./сут. Для коротких участков, резко отличающихся
от смежных (мосты, перекрёстки), длина из формулы исключается и
показатель измеряется в ДТП на миллион прошедших автомобилей
(формула~(1.2)). Для надёжной оценки документ требует данных
не менее чем за 3"--~5~лет.

Произведение $365 \cdot L \cdot N$ в формуле~\eqref{eq:rel_accident_rate}
есть годовой пробег по участку. Показатель $\mathit{И}$ выражает
число ДТП на единицу этого пробега.
Критерий отнесения участка к местам концентрации ДТП в Рекомендациях
ветвится в зависимости от того, известна ли интенсивность. При
$N > 3000$~авт./сут (п.~1.2.6) требуется одновременно, чтобы абсолютное
число ДТП за три года было не меньше табличного порога, зависящего от
интенсивности и длины участка, и чтобы $\mathit{И} \geq 0{,}3$. При
$N < 3000$~авт./сут (п.~1.2.9) используется критическая плотность ДТП
на километр в год, с отдельными порогами для дорог вне и внутри
населённых пунктов. Наконец, п.~1.2.10 предусматривает случай, когда
сведений об интенсивности нет вовсе: тогда временно, до получения
сведений, допускается оценивать участок по фактической плотности ДТП,
т.\,е. по числу происшествий на единицу длины.

Выявленные участки концентрации делятся по величине~$\mathit{И}$ на
малоопасные, опасные и очень опасные (п.~1.2.14, таблица~1.5), причём
пороги зависят от типа дороги: для двухполосных дорог по километровым
участкам они составляют соответственно 0{,}20"--~0{,}70, 0{,}70"--~1{,}30
и более 1{,}30~ДТП на миллион автомобиле"=километров, для многополосных
дорог с разделительной полосой "--- заметно ниже. Отдельно (п.~1.2.12)
участки классифицируются по динамике как прогрессирующие, стабильные и
регрессирующие "--- по статистической значимости изменения числа ДТП
относительно среднего за предыдущий период не менее трёх лет. Такая
проверка помогает отличать изменение от случайного колебания, но сама
по себе не устраняет смещение из-за отбора участков с необычно высоким
числом ДТП. Поправка на регрессию к среднему требует отдельной процедуры,
например эмпирического байесовского метода (подраздел~1.3).

\textbf{Метод коэффициентов аварийности.} Это наиболее известный
отечественный метод, обычно связываемый с именем Бабкова. Он не требует
статистики ДТП по конкретному участку и опирается на обобщённые
зависимости аварийности от дорожных условий. Итоговый коэффициент
аварийности участка вычисляется как произведение частных коэффициентов:
\begin{equation}
  K_{\text{итог}} = K_{1} \cdot K_{2} \cdot \ldots \cdot K_{18},
  \label{eq:babkov}
\end{equation}
где каждый $K_i$ "--- табличный множитель, характеризующий изменение
относительной аварийности под влиянием данного фактора по сравнению с
эталонным участком (прямая горизонтальная дорога с проезжей частью
7{,}5~м, укреплёнными обочинами и шероховатым покрытием). Восемнадцать
частных коэффициентов охватывают интенсивность движения ($K_1$), ширину
проезжей части и обочин, продольный уклон ($K_4$), радиус кривой в плане,
расстояние видимости, ширину мостов, длину прямых участков, тип и
видимость пересечений в одном уровне, число полос, расстояние до
застройки, длину населённого пункта и подходов к нему, коэффициент
сцепления покрытия ($K_{16}$), ширину разделительной полосы и близость
крутых откосов~\cite{babkov1993,vsn2586}.

По величине $K_{\text{итог}}$ участки на равнинной и пересечённой
местности классифицируются у Бабкова~\cite{babkov1993} как неопасные
(менее~10), малоопасные (10"--~20), опасные (20"--~40) и очень опасные
(более~40). Рекомендации Минтранса (пункты 1.4.7"--~1.4.11)
задают
практические пороги: участки с $K_{\text{итог}}$ выше 15"--~20 подлежат
перепроектированию при реконструкции, выше 25"--~40 (в холмистой
местности) "--- перестройке при капитальном ремонте, а на городских
улицах при реконструкции и новом строительстве не допускаются участки с
$K_{\text{итог}}$ выше~25.
Результат метода оформляется в виде линейного графика коэффициентов
аварийности: под сжатым планом и продольным профилем трассы строится
ступенчатая кривая $K_{\text{итог}}$ по километражу, пики которой и
есть опасные участки. При $K_{\text{итог}} > 15$ дополнительно
вычисляется стоимостной коэффициент, учитывающий тяжесть последствий
через произведение ещё четырнадцати частных коэффициентов тяжести.

Частные коэффициенты заданы таблицами для определённых дорожных
условий. Для городской сети существенны пешеходные потоки,
частые пересечения и плотность застройки, которые в этом методе
представлены обобщённо.

\textbf{Метод коэффициентов безопасности.} Оценивает не абсолютный
риск, а резкость изменения скоростного режима. Коэффициент безопасности
участка определяется как отношение максимальной скорости движения на
участке к максимальной скорости въезда на него (п.~1.3.1 Рекомендаций).
Чем сильнее участок заставляет водителя сбросить скорость по сравнению
с предшествующим, тем ниже коэффициент и тем опаснее участок; в
классификации Бабкова~\cite{babkov1993} участки с коэффициентом
ниже~0{,}4 считаются очень опасными, выше~0{,}8 "---
неопасными. В основе метода лежит представление о том, что риск связан
не с высокой или низкой скоростью участка как таковой, а с перепадами
скоростного режима. Эту идею можно соотнести с сетевыми признаками:
сменой класса улицы, выездом с магистрали на второстепенную сеть,
примыканием скоростного участка к регулируемому пересечению.

\textbf{Метод конфликтных ситуаций.} Основан на регистрации
изменений скорости и траектории движения автомобилем-лабораторией:
конфликтные ситуации делятся на лёгкие, средние и критические,
приводятся к критическим с весовыми коэффициентами, и по их числу на
миллион автомобиле"=километров участок относится к одной из степеней
опасности (§~1.6 Рекомендаций). С ним тесно связан применяемый для
пересечений \emph{метод конфликтных точек}~\cite{babkov1993}:
подсчитывается число точек отклонения, слияния и пересечения траекторий,
каждая со своим весом (пересечение существенно опаснее отклонения), и по
взвешенной сумме узел относится к простым, средней сложности, сложным
или очень сложным. Последний метод непосредственно связывает опасность с
топологией узла и допускает автоматизацию по данным о конфигурации
пересечения.

\textbf{ОДМ 218.6.015-2015 и понятие места концентрации ДТП.}
Методический документ Росавтодора содержит рекомендации по учёту и
анализу ДТП. В редакции распоряжения от 31 января 2017 года №~105-р
место концентрации ДТП определяется через пространственный масштаб
и повторяемость происшествий. Рассматривается участок дороги или
улицы длиной до 1000~м вне населённого пункта, до 200~м внутри него
либо пересечение дорог или улиц. За отчётный год на объекте должны
быть зарегистрированы не менее трёх ДТП одного вида либо не менее
пяти в сумме по всем видам; учитываются происшествия с погибшими
или ранеными~\cite{odm2015}. Данный критерий, в отличие от критериев
Рекомендаций Минтранса, не использует интенсивность движения и
опирается только на абсолютные числа происшествий, т.\,е. относится к
тому же классу, что и п.~1.2.10 Рекомендаций.

\subsection{Зарубежные методы: функции безопасности и системный подход}

Зарубежная методология количественной оценки безопасности дорог
наиболее полно изложена в руководстве \foreignlanguage{english}{Highway Safety Manual} (HSM)~\cite{hsm2010}, изданном в 2010~году Американской ассоциацией дорожных администраций штатов (AASHTO). Руководство не имеет силы закона,
но фактически является отраслевым стандартом и опирается на несколько
десятилетий исследований, обобщённых в работах Э.~Хауэра \cite{hauer1997}
и в обзоре Лорда и Маннеринга~\cite{lord2010}.

\textbf{Предсказательный метод.} Ожидаемое число ДТП на объекте
(перегоне или пересечении) за год представляется произведением трёх
сомножителей:
\begin{equation}
  N_{\text{predicted}} = N_{\text{SPF}} \cdot
  \prod_{i} \mathrm{CMF}_i \cdot C.
  \label{eq:hsm}
\end{equation}
Первый сомножитель "--- функция безопасности (Safety Performance
Function, SPF) "--- даёт базовое число ДТП для объекта данного типа при
стандартных условиях как функцию экспозиции. Например, для перегона
может использоваться лог-линейная спецификация
\begin{equation}
  N_{\text{SPF}} = \exp\!\left(a + b\log\mathrm{AADT} + c\log L\right),
  \label{eq:spf}
\end{equation}
где $\mathrm{AADT}$ и $L$ "--- численные значения среднегодовой
суточной интенсивности и длины в выбранных единицах (например,
авт./сут и км), а $a$, $b$, $c$ "--- параметры для данного типа
объектов и единиц измерения~\cite{lord2010,hsm2010}. Форма SPF
выбирается для конкретного типа объекта; для пересечений обычно
учитываются потоки главного и второстепенного направлений.
SPF и коэффициент относительной аварийности~\eqref{eq:rel_accident_rate}
связаны учётом экспозиции, но имеют разный смысл и размерность:
первая предсказывает число ДТП за период, второй выражает число
наблюдавшихся ДТП на единицу пробега.

Для функций безопасности часто используют отрицательную биномиальную
регрессию. В рассматриваемом ниже варианте эмпирического байесовского
метода дисперсия числа ДТП равна $\mu+k\mu^2$, где $\mu$ --- ожидаемое
число ДТП, а $k$ --- параметр сверхдисперсии.

Второй сомножитель "--- произведение поправочных коэффициентов (Crash
Modification Factor, CMF), по одному на каждое отклонение объекта от
базовых условий: более узкая полоса даёт $\mathrm{CMF} > 1$, наличие
освещения "--- $\mathrm{CMF} < 1$ и т.\,д. Часть~D руководства
представляет собой справочник таких коэффициентов, а Федеральное
управление автодорог США ведёт открытый каталог CMF Clearinghouse с
оценкой надёжности каждого значения. По назначению CMF близок к
частному коэффициенту $K_i$ из формулы Бабкова~\eqref{eq:babkov}:
оба отражают различия дорожных условий. Численные значения относятся
к базовым условиям и типам объектов, на которых они оценены.

Третий сомножитель "--- калибровочный коэффициент~$C$. Различия
между регионами и периодами могут смещать абсолютный уровень прогноза.
Один из способов локальной калибровки "--- отношение суммы наблюдённых
ДТП к сумме исходных прогнозов на отдельной местной выборке.
Существуют также калибровочные функции, зависящие от характеристик
объекта~\cite{fhwacalibrator2018}.
Постоянный положительный множитель корректирует общий уровень
прогноза при сохранении порядка объектов. Для его оценки нужна
местная история ДТП.

\textbf{Эмпирический байесовский метод.} Отбор участков по наблюдённому
числу ДТП страдает от регрессии к среднему (подраздел~1.1). HSM решает
эту проблему, комбинируя прогноз модели с наблюдением. Обозначим
через $S_{\text{predicted}}$ сумму годовых прогнозов за период
наблюдений, а через $S_{\text{observed}}$ "--- число фактических ДТП
за тот же период. Тогда оценка ожидаемого числа за этот период имеет вид
\begin{equation}
  S_{\text{expected}} = w S_{\text{predicted}}
  + (1-w)S_{\text{observed}}, \qquad
  w = \frac{1}{1+k S_{\text{predicted}}},
  \label{eq:eb}
\end{equation}
где $k$ "--- параметр сверхдисперсии функции безопасности, $w$ "---
вес прогноза, а $1-w$ "--- вес наблюдений. При фиксированных $k>0$
и положительном среднегодовом прогнозе увеличение периода наблюдений
снижает $w$. Среднегодовая оценка получается делением
$S_{\text{expected}}$ на число лет.

\textbf{Сетевой скрининг.} Часть~B руководства описывает процесс
управления безопасностью сети из шести шагов, первый из которых "---
сетевой скрининг (network screening): ранжирование всех объектов сети
по опасности и формирование короткого списка для последующей
диагностики. Диагностика выполняется инженером вручную по каждому
отобранному объекту, и именно это определяет требования к скринингу:
важнее не пропустить опасный объект, чем избежать лишнего в списке.
Руководство перечисляет более десятка показателей для ранжирования, от
простейших (частота ДТП, частота на миллион автомобиле"=километров) до
основанных на эмпирическом байесовском методе. Наиболее
содержательный из них "--- потенциал повышения безопасности (Potential
for Safety Improvement, PSI): разность между ожидаемым по
формуле~\eqref{eq:eb} числом ДТП и предсказанным функцией безопасности.
При согласованных периодах PSI характеризует превышение ожидаемой
аварийности над прогнозом для сопоставимых условий и интенсивности.
Такое превышение служит основанием для обследования участка и
поиска возможных причин повышенной аварийности.

\textbf{Системный подход.} Классический скрининг реактивен: он
выделяет места, где ДТП уже накопились. Системный подход к безопасности
(Systemic Safety Approach), оформленный Федеральным управлением автодорог
США в виде методики и программного инструмента~\cite{fhwa2013},
противопоставляет этому проактивную логику. Вместо поиска «горячих
точек» по статистике выявляются факторы риска "--- физические признаки
дороги, статистически связанные с тяжёлыми ДТП, "--- и мероприятия
разворачиваются на всех объектах сети, обладающих этими признаками,
включая те, где происшествий ещё не было.

\textbf{Инфраструктурные рейтинги.} Международная программа оценки
дорог iRAP~\cite{irap} использует инфраструктурные рейтинги. По видеосъёмке
кодируется более пятидесяти атрибутов дороги (разметка, число полос,
придорожные препятствия, устройство пересечений, тротуары), и для
каждого стометрового участка вычисляется рейтинг от одной до пяти
звёзд без обязательной местной истории ДТП для каждого оцениваемого
участка.

\textbf{Макроуровневые модели.} В академической литературе задача
оценки аварийности по территориальным единицам, а не по участкам дорог,
известна как макроуровневое моделирование (macro-level crash prediction).
ДТП агрегируются по зонам "--- обычно транспортным районам "--- и
моделируются отрицательной биномиальной регрессией по признакам зоны:
плотности сети, населению, землепользованию, долям дорог разных классов.
В более сложных моделях учитывают, что аварийность соседних зон может
быть связана. Один из таких подходов "--- модель
Бесага"--~Йорка"--~Молье~\cite{besag1991}. Обзор~\cite{lord2010}
систематизирует эти модели и их ограничения.

\subsection{Пространственная единица анализа: участок дороги или ячейка территории}

В дорожных методиках единицей анализа обычно служит однородный
перегон или перекрёсток. Для городской сети привязка ДТП к отдельному
ребру затруднена: на пересечении сходятся несколько улиц, а координаты
карточки могут быть смещены относительно места происшествия.
Кроме того, аварийность связана с окружением дороги --- застройкой,
остановками и объектами притяжения пешеходов.

При территориальном подходе события и дорожные характеристики
агрегируются по зонам. В макроуровневых моделях используются
транспортные районы и административные единицы. Их размеры и границы
влияют на результат; эта зависимость известна как проблема изменяемой
территориальной единицы (modifiable areal unit problem)~\cite{openshaw1984}.
Регулярная сетка задаёт воспроизводимое разбиение, масштаб которого
можно выбирать под задачу.

H3 --- иерархическая система геопространственной индексации
с преимущественно гексагональными ячейками~\cite{h3}. Гексагональная решётка удобна для описания
окружения: в идеальном плоском случае у ячейки шесть равноудалённых
соседей. В сферической H3 размеры ячеек и расстояния
различаются. Средняя площадь гексагона разрешения~9 составляет
около 0{,}105~км$^2$~\cite{h3stats}.

Размер ячейки выбирается между двумя требованиями: достаточно
событий для оценки и достаточно подробная локализация. Мелкая сетка
усиливает разреженность и чувствительность к смещениям координат;
крупная объединяет дороги с разными условиями движения. Разрешение~9
даёт масштаб в несколько сотен метров. С нормативным участком
длиной 200~м оно сопоставляется приближённо: одна ячейка может
содержать несколько дорог, а отдельный участок --- пересекать её границу.

Признаки ячейки описывают суммарную длину дорог по классам, устройство
пересечений, наличие регулирования, покрытие и рельеф. Результат
прогноза относится ко всей этой области, которую затем обследуют
подробнее. Агрегирование смягчает небольшие ошибки привязки,
но сохраняет чувствительность к переносу события через границу ячейки.

Время задаёт второе измерение панели: для каждой пары
«ячейка~$\times$~период» определяется число происшествий. Длительность
периода влияет на долю нулевых наблюдений и возможность учитывать
повторные ДТП.

\subsection{Открытые геоданные как источник характеристик дорожной сети}

Для построения признаков нужны геометрия дорог, сведения об их
устройстве и координаты происшествий. Рассматриваемые источники
различаются назначением и пространственным покрытием.

\textbf{Карточки ДТП.} Госавтоинспекция регистрирует дату и место
происшествия, его вид, дорожные условия, участников и тяжесть
последствий~\cite{odm2015}. Статистические показатели публикуются
на её портале~\cite{gibddstat}; машиночитаемые данные распространяет,
в частности, проект «Карта ДТП»~\cite{dtpstat}.

OpenStreetMap~\cite{osm} предоставляет геометрию дорожной сети и
теги объектов. По линиям дорог извлекаются класс, полосность,
покрытие, ограничение скорости и односторонность; по точкам ---
светофоры, переходы и другие элементы обустройства. Из самой геометрии
и общих узлов строятся показатели длины, связности и устройства
пересечений.

Полнота OSM зависит от территории и атрибута. У извлечённой дороги
есть геометрия и класс, тогда как полосность или покрытие могут
быть неизвестны. Для таких полей различаются пропуск и подтверждённое
значение: ненанесённый тег освещения, например, оставляет состояние
освещения неопределённым.

\textbf{Паспортные сведения СКДФ.} Федеральная государственная
информационная система «Система контроля за формированием и
использованием средств дорожных фондов»~\cite{skdf} содержит
категорию дороги, ширину проезжей части, покрытие, число полос
и для части объектов интенсивность движения. Данные доступны через
открытый картографический сервис WFS. Их сопоставление с OSM
требует согласовать геометрию и единицу учёта: паспорт может описывать
дорогу целиком, а OSM --- отдельные проезжие части. Поэтому полноту
оценивают по обеспеченной длине сети и отдельным атрибутам.

\textbf{Рельеф.} Copernicus DEM GLO-30~\cite{copernicusdem} ---
глобальная модель поверхности с разрешением около 30~м. По профилю
высот вдоль дороги оценивается продольный уклон, учитываемый и в
классических методиках ($K_4$ у Бабкова). Высоты считываются вдоль
дорог с шагом около 20~м. Уклон в каждой точке находят по наклону
прямой, подобранной к высотам в окне длиной около 200~м вдоль дороги.
Так уменьшается влияние локальных скачков высоты, в том числе
связанных со зданиями и растительностью. Для ячейки рассчитывается
средний модуль полученных уклонов с весами по длине дорожных отрезков.
Этот признак характеризует крутизну дорог в ячейке, а не разницу высот
между её границами.

Спутниковые композиты VIIRS~\cite{viirs} характеризуют ночную
радиояркость территории. Размер пикселя на широте Саратова
составляет около 460~$\times$~290~м. Такой показатель может отражать
освещённость и городскую активность, но усредняет разные объекты
внутри пикселя. Чувствительность прибора к светодиодному
освещению ограничена.

\subsection{Структурные факторы риска и проблема экспозиции}

Структурные признаки описывают медленно меняющиеся свойства места:
класс улицы, геометрию, пересечения, покрытие, уклон и окружение.
Погода и время суток характеризуют условия движения. Их влияние
может зависеть от места: снег по-разному действует на ровный участок
и дорогу с уклоном, а значение темноты связано с наличием освещения.

Экспозиция связывает частоту происшествий с объёмом движения.
На перегоне она может выражаться пробегом, на пересечении ---
входящими потоками. Для нормирования нужны сопоставимые измерения
по всей рассматриваемой сети. Скорость и пробочность описывают
состояние движения, а измерения на отдельных постах требуют модели
для распространения на остальную сеть.

При неполном учёте потока используют косвенные характеристики:
класс и полосность дороги, связность сети, плотность застройки
и объектов притяжения. В них смешиваются сведения о загруженности
и дорожных условиях. Например, магистраль может иметь много ДТП
из-за большого числа поездок при низкой частоте происшествий
на единицу пробега. По структурным признакам в таком случае
оценивается ожидаемое число ДТП; относительная аварийность
по формуле~\eqref{eq:rel_accident_rate} требует данных о пробеге.

Одним из косвенных показателей служит центральность по кратчайшим
путям. Проходящие через ячейку пути отражают её транзитную роль
при заданном выборе пунктов отправления и назначения и весов
рёбер. Для автомобильного движения важны время проезда,
односторонность и разрешённые манёвры.

\subsection{Машинное обучение для редких пространственных событий}

Число ДТП можно прогнозировать с помощью пуассоновской или
отрицательной биномиальной регрессии~\cite{lord2010}. В таких моделях
связь признаков с ожидаемым числом ДТП задаётся формулой. Если нужно
учесть, что влияние признака меняется в разных условиях, это также
приходится указывать в формуле. Деревья решений могут находить такие
зависимости при обучении на данных, поэтому они удобны при большом
числе дорожных и пространственных признаков.

Градиентный бустинг последовательно строит деревья, приближающие
направление уменьшения функции потерь~\cite{friedman2001}. Для прогноза
неотрицательного среднего используется логарифмическая шкала:
\begin{equation}
  a(\mathbf{x})=b_0+\sum_{m=1}^{M}\eta f_m(\mathbf{x}),\qquad
  \widehat\mu(\mathbf{x})=\exp a(\mathbf{x}),
  \label{eq:gbdt}
\end{equation}
где $b_0$ --- начальная константа, $\eta$ --- шаг обучения, а $f_m$ ---
деревья. Внутри дерева могут сочетаться несколько признаков,
что позволяет учитывать их взаимодействия. Одна из реализаций
бустинга --- CatBoost~\cite{catboost2018}.

Пусть $y$ --- число ДТП в ячейке за год, а $a$ --- выход модели
до преобразования в ожидаемое число ДТП $\exp(a)$.
Пуассоновские потери, с точностью до константы, не зависящей
от прогноза, имеют вид
\begin{equation}
  L_{\mathrm{P}}(y,a)=\exp(a)-ya.
  \label{eq:poisson_loss}
\end{equation}
Другой критерий обучения --- потери Tweedie:
\begin{equation}
  L_{\mathrm{T}}(y,a;p)=
  \frac{\exp((2-p)a)}{2-p}
  -y\frac{\exp((1-p)a)}{1-p},\qquad 1<p<2.
  \label{eq:tweedie_loss}
\end{equation}
Их реализация описана в документации CatBoost~\cite{catboostlosses}.
Параметр $p$ задаёт форму потерь Tweedie.

При бинарной цели положительной считается ячейка с хотя бы одним
ДТП. В этом случае применяется логистическая функция потерь,
а ансамбль оценивает логарифм шансов. Счётная цель дополнительно
сохраняет информацию о повторных происшествиях, что существенно
при отборе мест с большим ожидаемым числом ДТП.

Выбор признаков опирается на заполненность полей и сравнение
вариантов модели. Для проверки вклада признака его исключают
и повторяют оценку на той же выборке. Глобальные важности
и локальные вклады описывают участие признаков в прогнозе.
Причинные эффекты переустройства дороги изучаются отдельно,
например в исследованиях «до и после»~\cite{hauer1997}.

При наличии истории в модель включают прошлые числа ДТП или
используют их совместно со структурным прогнозом по
формуле~\eqref{eq:eb}. Для нового города без местной истории
структурная модель остаётся доступным вариантом. При сравнении
таких методов фиксируются состав данных, горизонт прогноза
и доступная каждому методу информация.

\subsection{Оценка качества моделей ранжирования риска}

Пусть в городе $n$ дорожных ячеек. Для ячейки $i$ обозначим
через $y_i$ фактическое число ДТП за проверяемый период, а через
$\widehat y_i$ --- прогноз модели. Ячейки сортируют по убыванию
прогноза и выбирают первые $k$. Для coverage@10 это верхние
10\,\% дорожных ячеек; например, из 1000 ячеек выбирают 100.
В общем случае $k$ равно доле $q$ от числа ячеек с округлением
вниз, но не менее одной ячейки. Обозначим выбранный набор через
$S_q$. Если за период было хотя бы одно ДТП, доля ДТП
в выбранных ячейках равна
\begin{equation}
  \operatorname{Coverage}@q=
  \frac{\sum_{i\in S_q}y_i}{\sum_{i=1}^{n}y_i}.
  \label{eq:coverage}
\end{equation}
Например, если во всём городе произошло 300 ДТП, а в выбранных
100 ячейках --- 120, то coverage@10 равно 120/300, или 40\,\%.
Ячейки выбираются по прогнозу, а число ДТП подсчитывается уже
после отбора. Граница отбора проходит между $k$-й и следующей
ячейками в списке. Если их прогнозы равны, для coverage берут
средний результат по возможным вариантам выбора среди ячеек
с одинаковым прогнозом. Так порядок строк в данных не влияет
на результат.

ROC-AUC и средняя точность (average precision, AP) показывают,
насколько высоко модель ставит ячейки с ДТП по сравнению с ячейками
без ДТП~\cite{saito2015}. AP объединяет значения точности при разных
порогах отбора~\cite{sklearnap}. Корреляция Спирмена показывает,
насколько порядок прогнозов согласуется с числом ДТП.

Для оценки суммарной калибровки сравнивают сумму прогнозов
с фактическим числом ДТП за тот же период:
\begin{equation}
  A=\frac{\sum_i\widehat y_i}{\sum_i y_i}.
  \label{eq:calibration}
\end{equation}
Значение $A$ около единицы означает, что общий прогноз близок
к фактическому числу ДТП. Ошибки между ячейками дополнительно
сравнивают по группам прогнозируемого уровня и с помощью
пуассоновского девианса.

Способ проверки зависит от задачи. Если прогнозируются будущие годы
в знакомом городе, модель обучают на более ранних годах и проверяют
на более поздних. Для проверки переноса на новый город его целиком
исключают из обучения (leave-one-city-out, LOCO). Это показывает,
как модель работает при других дорожных признаках и качестве
исходной карты~\cite{roberts2017}.

При сравнении городов метрика сначала вычисляется внутри каждого
из них, а затем усредняется с равными весами:
\begin{equation}
  \operatorname{MacroCoverage}@q=
  \frac{1}{G}\sum_{g=1}^{G}\operatorname{Coverage}_g@q,
  \label{eq:macro}
\end{equation}
где $G$ --- число проверяемых городов. Кроме среднего, приводят
нижний квартиль и худший результат. Если ДТП распределены по городу
более равномерно, небольшая доля ячеек охватывает меньше событий.

При обучении город с большим числом строк сильнее влияет на модель.
Чтобы суммарный вес каждого города был одинаковым, каждой его строке
назначают вес
\begin{equation}
  w_g=\frac{N}{G\,n_g},
  \label{eq:city_weights}
\end{equation}
где $N$ --- общее число обучающих строк, $G$ --- число обучающих
городов, а $n_g$ --- число строк города $g$. Средний вес строки
при этом равен единице.

Обучение повторяют с разными состояниями генератора случайных чисел,
чтобы оценить разброс результатов. При сравнении двух моделей
используют парный бутстреп: для обеих моделей берут одни и те же
повторные выборки ячеек. Если соседние ячейки зависимы,
в выборку можно брать целые пространственные блоки.

\textbf{Предел предсказуемости.} Сортировка по будущему числу ДТП
использует информацию, которой на момент прогноза нет.
Для более реалистичного ориентира каждой ячейке задают скрытую
среднегодовую частоту ДТП $\lambda_i$ и симулируют число событий
за $H$ лет:
\begin{equation}
  Y_i\mid\lambda_i\sim\operatorname{Poisson}(H\lambda_i).
  \label{eq:oracle_poisson}
\end{equation}
Оракул выбирает ячейки по $\lambda_i$, а его результат оценивается
по симулированным $Y_i$. Если частоты $\lambda_i$ различаются по
гамма-распределению, число ДТП имеет отрицательное биномиальное
распределение. Повторные симуляции учитывают неопределённость
самих частот.

% ---------- 2 Постановка практической задачи ----------

\section{Постановка практической задачи}

\textbf{Основная территория и данные.} Основное исследование
проводится на саратовской территории, ограниченной прямоугольником
51{,}40--51{,}70$^\circ$~с.\,ш. и
45{,}80--46{,}20$^\circ$~в.\,д. Это исследовательская область,
которая включает часть Энгельса и не совпадает с административной
границей Саратова. Из сетки H3 разрешения~9 отбираются ячейки
с ненулевой длиной автомобильных дорог OSM: 4\,296 ячеек.
Период наблюдения --- 2015--2025 годы. Карточки ДТП с пострадавшими
относятся к ячейкам по координатам; в основной дорожной панели
учитываются только события выбранных дорожных ячеек.

Исходная выгрузка получена 01.07.2026 с сайта проекта
«Карта ДТП»~\cite{dtpstat} и содержит около 1{,}6~миллиона карточек
по стране с 2015 года. Первоисточник сведений --- Госавтоинспекция.
На странице выгрузки проект сообщает о внесении изменений в данные,
включая корректировки координат. Последующая обработка ведётся
по сохранённому локальному файлу.

\textbf{Развитие постановки.} Исходный этап использовал 132 месяца
и бинарную цель --- наличие хотя бы одного ДТП в следующем месяце.
Он служил для построения базовой модели и последовательной проверки
групп признаков. Основная постановка перешла к годовому числу ДТП,
чтобы учитывать повторные происшествия и сопоставлять структурный
прогноз с историей, EB и критериями концентрации.

Годовая панель содержит 47\,256 наблюдений «ячейка~$\times$~год»
и 15\,927 ДТП за 2015--2025 годы. Для ячейки $i$ и года $t$
определяются:
\begin{itemize}
  \item целевая переменная $Y_{i,t}$ --- число ДТП с пострадавшими;
  \item структурные признаки $\mathbf{x}_i$ --- характеристики сети
        OSM, геометрии и устройства пересечений, атрибутов дорог,
        рельефа и паспортных сведений СКДФ в зависимости от варианта;
  \item дополнительные исторические признаки --- числа ДТП за
        предшествующие периоды, используемые только в вариантах,
        где местная история доступна.
\end{itemize}
Геометрия современного среза OSM используется для расчёта
статичных дорожных признаков на период 2015--2025 годов.

Для CatBoost числовые признаки ограничиваются обучающим
99-м процентилем. Пороги вычисляются заново в каждом разбиении,
до применения модели к проверочным наблюдениям.

Измеренная интенсивность движения в набор признаков не входит.

\textbf{Задача и схема проверки.} Требуется построить прогноз
ожидаемого годового числа ДТП
\begin{equation}
  \widehat\mu_i=f(\mathbf{x}_i)
\end{equation}
и ранжировать по нему дорожные ячейки. Выбор функции потерь и
вариантов выполняется на валидационных годах 2021 и 2022:
обучение соответственно на 2015--2020 и 2015--2021 годах.
Основная проверка проводится на 2023--2025 годах после обучения
на 2015--2022 годах. Прогноз зафиксирован до начала проверочного
периода; для оценки суммы ДТП за три года используется
$\widehat Y_i=3\widehat\mu_i$. В сценарии с историей признаки также
строятся только по сведениям, доступным к моменту прогноза.

Качество определяется coverage@1, coverage@5 и coverage@10
по формуле~\eqref{eq:coverage}, корреляцией Спирмена и показателями
калибровки. ROC-AUC и AP по наличию хотя бы одного события
используются как вспомогательные метрики. Прямые сравнения вариантов
проводятся на одних и тех же ячейках и периодах.

\textbf{Модели и варианты.} Начальная месячная базовая модель
CatBoost использовала 61 структурный признак, балансировку классов
и ограничение признаков по обучающему 99-му процентилю. Её метрики
относятся только к этому этапу. После последовательных проверок
вклада признаков саратовский структурный вариант содержит 32 признака.
Для годовой постановки сопоставлены бинарный отклик, Poisson и
Tweedie. Выбранный вариант --- CatBoost с потерями Tweedie при
$p=1{,}3$, 800 итерациями, шагом 0{,}05 и глубиной~6;
балансировка бинарных классов в нём не применяется.
Отдельно сравниваются ранжирование по истории, совместная модель
«структура и история» и EB, в том числе с перекрёстным обучением
структурной модели по пространственным блокам. В последнем случае
прогноз для обучающей ячейки получается без её группы блоков.
Простым структурным ориентиром служит суммарная длина дорог старших
классов внутри ячейки.

В месячных опытах дополнительно проверялись календарные признаки
и ночная яркость VIIRS. Устойчивого улучшения основных метрик
ранжирования они не дали; основной годовой вариант оставлен структурным.
Численные результаты этих сравнений будут представлены
в практической части.

\textbf{Межгородская проверка.} Подход дополнительно
проверен на 14 городах с населением более 350 тысяч человек из
локальной выгрузки OSM Поволжья: Казань, Нижний Новгород, Уфа,
Самара, Пермь, Саратов, Тольятти, Ижевск, Ульяновск, Набережные
Челны, Оренбург, Чебоксары, Пенза и Киров. Для порогового отбора
использованы оценки населения на 2025 год из инвентаря городов:
теги OSM с указанием года и источника, а для Кирова --- сведения
региональной статистики. Население служит критерием отбора городов.
Территория анализа задаётся муниципальными полигонами из снимка OSM
от 3 июня 2026 года, включая входящую в них сельскую периферию.
Ячейка включается в когорту по
положению её центра внутри полигона; для анализа чувствительности
отдельно оценивается подвыборка в радиусе 15~км от центра города.

Межгородской набор содержит 28 OSM-признаков, без СКДФ и DEM. В каждой проверке с исключением одного города
целевой город полностью исключён из обучения и вычисления 99-го
процентиля признаков. Модель обучается на остальных
13 городах, а прогноз 2023--2025 годов оценивается в исключённом
городе. Основные окна обучения --- 2015--2022 и 2017--2022 годы;
они сопоставляются для проверки варианта без данных за 2015--2016 годы.
Дополнительно
сравниваются равные веса строк и равный суммарный вес обучающих
городов по формуле~\eqref{eq:city_weights}.

Качество сначала оценивается внутри каждого города, затем по
равновесному среднему~\eqref{eq:macro}, нижнему квартилю, худшему
результату и калибровке. Результаты по 14 городам используются
для анализа межгородских различий и формирования гипотез.

\textbf{Исследовательские вопросы.} При анализе результатов
рассматриваются следующие вопросы:
\begin{enumerate}
  \item какие группы структурных признаков дают устойчивый вклад
        в качество на одинаковой когорте;
  \item как смена месячной бинарной цели на годовой счёт влияет
        на ранжирование и калибровку;
  \item как структура соотносится с историей ДТП и их комбинациями
        при разной доступности местной информации;
  \item насколько качество различается между городами и как это
        связано с распределениями признаков, границами территорий
        и полнотой картографических данных;
  \item меняются ли выводы при исключении ранних лет и при равном
        вкладе городов в обучение;
  \item как результаты соотносятся с приближённым ячеечным критерием
        концентрации по ОДМ и с условным пределом предсказуемости
        в симуляционной модели.
\end{enumerate}

\textbf{Оценка устойчивости.} Модели обучаются с несколькими
начальными состояниями генератора. В саратовском сравнении
разность метрик оценивается парным бутстрепом по ячейкам
(1000 повторных выборок). В межгородских сравнениях окон и весов
используются 500 повторных выборок пространственных блоков
H3 разрешения~7, чтобы учитывать локальную зависимость.

Согласованность координат с дорожной сетью изучается по событиям,
попавшим вне дорожных ячеек; их распределение сравнивается по годам.

Для диагностической симуляции параметры скрытого риска оцениваются
по истории 2015--2025 годов, для прогнозных сравнений ---
по данным до начала проверочного периода.
Параметры подготовки, начальные состояния и результаты
сохраняются в протоколе экспериментов.

% ---------- 3 Практическая часть (заголовки предварительные) ----------
\section{Данные и их подготовка}
% Временный сводный пункт: заменить развёрнутым описанием практики.
\subsection*{Что уже выполнено}

\begin{itemize}
  \item Подготовлена годовая панель ДТП и дорожных ячеек H3 для
        саратовской территории: 4\,296 ячеек и 47\,256 наблюдений
        за 2015--2025 годы.
  \item Из OSM получены признаки сети, геометрии и пересечений;
        отдельно подготовлены варианты с рельефом и сведениями СКДФ.
  \item Обучены месячные и годовые модели CatBoost. Годовой
        структурный вариант с 32 признаками сопоставлен с историей
        ДТП, её сочетанием с дорожными признаками и методом EB.
  \item Проведены проверки по времени и переноса на 14 городов;
        сравнены периоды обучения и способы взвешивания городов.
\end{itemize}

\textbf{Предварительные итоги.} В саратовском годовом сравнении
выбран CatBoost с потерями Tweedie; различия между счётными
вариантами по охвату ДТП невелики. При переносе между 14 городами
coverage@10 на проверке 2023--2025 годов составил от 47 до 77\,\%
в зависимости от города, в среднем около 61\,\%. Из-за подозрений
на ошибки координат в ранней части выгрузки отдельно проверено
обучение без данных за 2015--2016 годы. Этот вариант и равные
суммарные веса городов не дали устойчивого улучшения во всех городах.

% TODO: источники, сетка H3, построение признаков

\section{Разведочный анализ признаков}
% TODO

\section{Модели и результаты}
% TODO: месячный бейзлайн и абляции; годовые модели; история и EB; перенос между городами

\section{Сопоставление с классическими методами и ограничения}
% TODO

% ---------- Заключение ----------
\conclusion
% TODO

\bibliographystyle{ugost2003}
\bibliography{thesis}

\appendix

\end{document}
